\section{Order, gaps, and a scheduling bound}\label{sec:gaps}

So far we have watched knots one at a time. This section watches the
\emph{spaces between them}, and the payoff is a second proof that Pisot
parameters have bounded knot counts -- shorter than the automaton argument of
Section~\ref{sec:pisot}, and effective: it converts a single number, the
smallest gap in the orbit of $1/2$, directly into a bound on
$N_\lambda$.

Throughout, fix a run and consider the knots alive at a given moment, listed in
increasing order of position. Three elementary lemmas control everything.

\begin{lemma}[order]\label{lem:order}
Every move preserves the spatial order of the knots that survive it.
\end{lemma}

\begin{proof}
The maps $f_0$ and $f_1$ are increasing, so moves $L$ and $R$, which apply one
map to every survivor, preserve order. At a move $M$ the survivors below the
deleted interval are carried by $f_0$ into $(0,1/2)$ and those above it by
$f_1$ into $(1/2,1)$; each group keeps its internal order, and the first group
lands entirely to the left of the second.
\end{proof}

Thus two knots, once created, keep their left-to-right relationship for as long
as both live. Say that a pair of knots \emph{straddles} a move $M$ if one lies
below the deleted interval and the other above it.

\begin{lemma}[gap law]\label{lem:gaplaw}
Let $x<y$ be knots that both survive a move. If the move is $L$, is $R$, or is
an $M$ that the pair does not straddle, then the new distance is
$\lambda(y-x)$. If the pair straddles an $M$, the new distance is
$\lambda(y-x)-(\lambda-1)$.
\end{lemma}

\begin{proof}
In the first case both knots are carried by the same affine map of slope
$\lambda$. In the second, $x\mapsto\lambda x$ and
$y\mapsto\lambda y-(\lambda-1)$.
\end{proof}

So every distance between coexisting knots is stretched by the factor
$\lambda$ at every single move, with one exception: a pair that straddles an
$M$ has its distance reduced by $\lambda-1$ relative to that stretch. Since
all knots lie in $(0,1)$, no distance can ever reach $1$; the stretching must
therefore be paid down, and straddling is the only currency. The final lemma
says the currency is scarce.

\begin{lemma}[one service per move]\label{lem:service}
At any move, among the spatially consecutive pairs of any fixed set of
surviving knots, at most one pair straddles.
\end{lemma}

\begin{proof}
A straddling pair has its two members on opposite sides of the deleted
interval, so a consecutive pair straddles exactly when the deleted interval
falls between its two members -- and an interval falls between the members of
at most one consecutive pair of an ordered list.
\end{proof}

These three facts assemble into a pigeonhole argument.

\begin{theorem}[scheduling bound]\label{thm:schedule}
Suppose a run reaches $k$ simultaneous knots, and that throughout the run any
two coexisting knots are at distance at least $\delta>0$. Put
$W=\lceil\log_\lambda(1/\delta)\rceil$. Then
\[
  k\ \leq\ W+\lceil W/2\rceil+1.
\]
\end{theorem}

\begin{proof}
Fix a run with $k$ knots alive at time $n$, and look at the window consisting
of the last $W$ moves. Call a knot \emph{old} if its age is at least $W$; old
knots were alive throughout the window. By Lemma~\ref{lem:order} the old
knots keep a fixed spatial order, so their consecutive pairs are well defined
throughout. Consider one such pair. At time $n-W$ its distance was at least
$\delta$; by Lemma~\ref{lem:gaplaw} that distance is multiplied by $\lambda$
at each of the $W$ moves of the window unless the pair straddles one of them;
and $\lambda^W\delta\geq1$ while distances stay below $1$. So each
consecutive old pair straddles at least one move of the window. By
Lemma~\ref{lem:service} each move is straddled by at most one such pair.
Hence (number of old knots) $-\,1\leq W$.

The remaining knots have age less than $W$, and by the argument of
Proposition~\ref{prop:lowerbound} their birth times are at least two apart, so
there are at most $\lceil W/2\rceil$ of them. Adding the two counts gives the
bound.
\end{proof}

\begin{corollary}[Pisot boundedness, again, and effectively]\label{cor:effective}
Let $\lambda$ be Pisot, let $O$ be the (finite) orbit of $1/2$, and let
$\delta_0$ be the smallest distance between two points of $O$. Then
$N_\lambda\leq W+\lceil W/2\rceil+1$ with
$W=\lceil\log_\lambda(1/\delta_0)\rceil$.
\end{corollary}

\begin{proof}
Every knot is born at $1/2$ and moved by $f_0$ and $f_1$, so all positions lie
in $O$; coexisting knots are distinct by Lemma~\ref{lem:distinct}; so any two
are at distance at least $\delta_0$, and Theorem~\ref{thm:schedule} applies.
\end{proof}

For the golden ratio the orbit is the five-point set of
Theorem~\ref{thm:phi}, whose smallest gap is
$\frac{\varphi-1}{2}-\frac{2-\varphi}{2}=\varphi-\tfrac32=(\sqrt5-2)/2\approx
0.118$; then $W=5$ and $N_\varphi\leq9$. For the plastic number the orbit has
$153$ points and smallest gap $\approx0.00239$ (a $50$-digit computation over
the orbit), giving $W=22$ and $N_\rho\leq34$. The true values are $2$ and $7$,
so the bound is loose by a factor of about five -- but it needs no reachable
configurations and no automaton, only the orbit and one minimum, and unlike
Theorem~\ref{thm:pisot} it is uniform: the same three lemmas serve every
parameter.

Read in the other direction the theorem says that large knot counts
\emph{force near-collisions}.

\begin{corollary}[near-collisions are necessary]\label{cor:nearcollision}
If some run reaches $k$ simultaneous knots, then for every integer
$m<2(k-2)/3$ there is a moment at which two coexisting knots lie within
$\lambda^{-m}$ of one another.
\end{corollary}

\begin{proof}
If every pair of coexisting knots always kept distance at least
$\lambda^{-m}$, the theorem with $\delta=\lambda^{-m}$, $W=m$ would give
$k\leq m+\lceil m/2\rceil+1\leq \tfrac32m+2$.
\end{proof}

At $\lambda=3/2$ this connects the knot count to the arithmetic of
Section~\ref{sec:threehalves}: positions of knots of age at most $a$ are odd
multiples of $2^{-(a+1)}$, so a pair within $\lambda^{-m}$ is a
near-coincidence among the integers $3^a-\sum\varepsilon_j2^j3^{a-j}$, and how
closely those subset sums can approach one another governs how many knots can
coexist. The erratic behaviour of such approximations is, we believe, the
arithmetic face of the erratic growth of $d_{3/2}(k)$.

Finally, the gap law suggests an accounting heuristic for the quadratic growth
observed at $\lambda=(1+\sqrt2)/2$, which we record without claiming a proof:
with $k$ knots there are $k-1$ consecutive distances, every one of them
stretched by $\lambda$ at every move, and Lemma~\ref{lem:service} allows only
one reduction per move. Each distance must be serviced before it grows to
length $1$, so the servicing load grows linearly with $k$ while the service
rate is fixed; sustaining $k$ knots therefore forces the configuration to
spend an ever-larger share of its moves on maintenance, and the cost of adding
one more knot grows roughly linearly in $k$ -- which integrates to
$d(k)\asymp k^2$. Making this rigorous would require a lower bound on how
small the serviced distances can be made, which is exactly the
anti-concentration question left open above.

\section{Periodic runs and Littlewood polynomials}\label{sec:kind}
